\section{One-sided boundary-logarithm approximation through zeros}\label{sec:boundary-log}

The quadratic floor concerns a sufficient positive surrogate. We now
retain the measured logarithm itself. The local theorem below controls
how much an absolute field error can raise its integral, even at
multiple real zeros. Its starting data are a nonzero anchor away
from the real interval and an upper bound on a larger disk.

\subsection{A one-dimensional sublevel theorem}

Fix the nested radii and constants
\begin{equation}\label{eq:disk-constants}
 R_0=\frac92,\quad \rho_0=\frac{13}3,\quad r_0=\frac{25}6,\quad
 B=R_0+r_0=\frac{26}3,\quad
 \kappa=\log\frac{27}{26},\quad
 \frac{\rho_0+r_0}{\rho_0-r_0}=51 .
\end{equation}

\begin{theorem}[Line sublevels and a logarithmic upper envelope]\label{thm:local-log}
Let \(f\) be holomorphic on a neighborhood of the closed disk
\(\overline{D(0,R_0)}\), with \(|f|\le1\) there and
\(|f(0)|\ge m>0\). Put
\begin{equation}\label{eq:n-C}
 A=\log(1/m),\qquad
 n=\max\{2,\lceil A/\kappa\rceil\},\qquad
 C_f=2Bn\ee^{51A/n}.
\end{equation}
Let \(\Gamma\) be a horizontal segment of length at most one inside
\(\overline{D(0,r_0)}\). Then, for every \(t>0\),
\begin{equation}\label{eq:sublevel}
 |\{z\in\Gamma:|f(z)|<t\}|\le C_ft^{1/n}.
\end{equation}
If \(g\) is measurable on \(\Gamma\) and
\(|f-g|\le\varepsilon\), with \(\varepsilon>0\), then
\begin{equation}\label{eq:local-envelope}
 0\le\int_\Gamma
       [\log(|g|+\varepsilon)-\log|f|]\dd\ell
 \le\frac{C_fn^2}{n-1}(2\varepsilon)^{1/n}.
\end{equation}
No lower bound for \(|f|\) on \(\Gamma\), simplicity or separation
of its zeros, or nonvanishing or holomorphy of \(g\), is assumed.
\end{theorem}

\begin{proof}
Jensen's formula \cite{Tao2020LocalZeros} bounds the number \(k\)
of zeros inside the open \(\rho_0\)-disk, including multiplicity,
by \(k\kappa\le A\). Remove exactly these zeros with disk Blaschke
factors:
\[
 b_j(z)=\frac{R_0(z-a_j)}{R_0^2-\overline{a_j}z},\qquad
 h(z)=f(z)\Big/\prod_{j=1}^k b_j(z).
\]
Fill the removable singularities. Since every factor has boundary
modulus one, the maximum principle gives \(|h|\le1\) on the outer
disk. It is zero-free on the open \(\rho_0\)-disk and
\(|h(0)|\ge m\). Harnack's inequality for the nonnegative harmonic
function \(-\log|h|\) gives
\(|h(z)|\ge\ee^{-51A}\) when \(|z|\le r_0\)
\cite{Tao2016Harnack}. If \(h\) has a zero on the \(\rho_0\)-circle,
apply Harnack on smaller disks and pass to the limit.
The elementary denominator bound for each Blaschke factor gives
\begin{equation}\label{eq:factor-lower}
 |f(z)|\ge\ee^{-51A}
       \prod_{j=1}^k\frac{|z-a_j|}{B}\qquad(|z|\le r_0).
\end{equation}
Every distance ratio is less than one.

Let \(q=\ee^{51A}t\). If \(q<1\) and \(k=0\), the sublevel set
is empty. Otherwise, when \(|f(z)|<t\), at least one distance
is below \(Bq^{1/k}\). The set is covered by \(k\) disks of this
radius. A disk of radius \(b\) meets a line in length at most \(2b\),
so the total length is at most
\[
 2Bkq^{1/k}\le2Bnq^{1/n}=C_ft^{1/n}.
\]
For \(q\ge1\), this bound exceeds the segment length and is
automatic. This proves \eqref{eq:sublevel} in one dimension.
An area estimate would not suffice for this trace argument.

The triangle inequality gives, off the finitely many segment zeros,
\[
 0\le\log(|g|+\varepsilon)-\log|f|
 \le\log(1+2\varepsilon/|f|).
\]
Write \(h_*=2\varepsilon\). Tonelli and the layer-cake identity give
\[
 \int_\Gamma\log(1+h_*/|f|)\dd\ell
 =\int_0^\infty
 \frac{h_*}{t(t+h_*)}|\{|f|<t\}\cap\Gamma|\dd t.
\]
On \(0<t<h_*\), bound the kernel by \(1/t\); on \(t>h_*\),
bound it by \(h_*/t^2\). Equation~\eqref{eq:sublevel} then gives
the two bounds \(C_fn h_*^{1/n}\) and
\(C_fn h_*^{1/n}/(n-1)\). Their sum is
\eqref{eq:local-envelope}. The same sublevel integral proves local
integrability of \(\log^-|f|\). The choice \(n\ge2\) makes the
second integral finite and avoids the endpoint logarithm at \(n=1\).
\end{proof}

An immediate consequence, for \(M\ge1\), is
\begin{equation}\label{eq:absolute-log-local}
 \int_\Gamma|\log|Mf||\dd\ell\le\log M+C_fn.
\end{equation}
Indeed \(|\log|Mf||\le\log M+\log^-|f|\), and integrating
\eqref{eq:sublevel} over \(0<t<1\) with measure \(\dd t/t\)
bounds the latter integral by \(C_fn\).

\subsection{The canonical anchor and explicit height dependence}

For \(v\in\R\), center these disks at \(c_v=v+4\ii\), and take
the segment \(I_v=[v-1/2,v+1/2]\). Its maximum distance from
the center is \(\sqrt{65}/2<r_0\).
Throughout the outer disk, \(1/4\le\Re s\le19/4\).
The pole-cancelled fractional-part identity
\eqref{eq:fractional} gives
\[
 |F(z)|\le2|v|+25.
\]
In detail \(|s-1|\le |v|/2+6\), while the integral in that
identity is bounded by \(1/\Re s\le4\). At the center,
the absolutely convergent Euler product gives
\begin{equation}\label{eq:canonical-anchor}
 |F(c_v)|\ge c_0:=\frac35\frac{\zeta(5)}{\zeta(5/2)}>0.
\end{equation}
This safe-line estimate is unconditional and supplies no lower
modulus on the real segment.

Define
\begin{equation}\label{eq:local-canonical-constants}
 M_v=2|v|+26,\quad A_v=\log(M_v/c_0),\quad
 n_v=\max\{2,\lceil A_v/\kappa\rceil\},\quad
 C_v=2Bn_v\ee^{51A_v/n_v}.
\end{equation}
Apply Theorem~\ref{thm:local-log} to \(F/M_v\).

\begin{corollary}[Canonical local acquisition]\label{cor:canonical-log}
If \(\sup_{I_v}|F-\widehat F|\le\delta\), with \(\delta>0\), then
\begin{equation}\label{eq:canonical-envelope}
 0\le\int_{I_v}
 [\log(|\widehat F|+\delta)-\log|F|]\dd x
 \le\mathcal E_v(\delta)
 :=\frac{C_vn_v^2}{n_v-1}
       \left(\frac{2\delta}{M_v}\right)^{1/n_v}.
\end{equation}
Also
\begin{equation}\label{eq:canonical-absolute-log}
 \int_{I_v}|\log|F(x)||\dd x\le\log M_v+C_vn_v.
\end{equation}
\end{corollary}
\begin{proof}
Use \(f=F/M_v\), \(g=\widehat F/M_v\),
\(\varepsilon=\delta/M_v\). Logarithmic constants cancel
exactly in the difference. Equation~\eqref{eq:absolute-log-local}
gives the second assertion.
\end{proof}

Because \(A_v/n_v\le\kappa\), the prefactor in
\eqref{eq:canonical-envelope} is \(O(\log^2(2+|v|))\),
whereas \(n_v\) is of order \(\log(2+|v|)\). To make the local
error at most \(h>0\), it suffices to take
\begin{equation}\label{eq:precision}
 \delta\le\frac{M_v}{2}
 \left(\frac{h(n_v-1)}{C_vn_v^2}\right)^{n_v}.
\end{equation}
This is finite precision for each fixed interval, not a uniform
positive Hölder exponent or a useful complexity theorem.

The local double-zero model illustrates the relevant scale without
any numerical approximation:
\begin{equation}\label{eq:double-log}
 \int_\R[\log(x^2+\delta)-\log x^2]\dd x=2\pi\sqrt\delta.
\end{equation}
Substitute \(x=\sqrt\delta\,u\) and integrate
\(\log(1+u^{-2})\) by parts on the positive half-line; its integral
is \(2\int_0^\infty(1+u^2)^{-1}\dd u=\pi\).
Thus a real zero produces no fixed error floor, although a Lipschitz
error bound fails even in this simple model. This is an algebraic
scale calculation, not a canonical collision control, and the
conservative exponent in \eqref{eq:canonical-envelope} is not
claimed sharp.

The approximant in this theorem may have zeros, and its logarithmic
integral is retained. This is the precise distinction from a
zero-free harmonic-baseline replacement. Likewise no small
quadratic-residual hypothesis from Section~\ref{sec:quadratic}
is used. The theorem measures signed information already present
in the exact field; it does not establish its favorable value.

